\documentclass[10pt,a4paper]{amsart}
\usepackage[margin=19mm]{geometry}
\usepackage{amsmath,amssymb,mathtools}
\usepackage[scr=rsfs,cal=euler]{mathalfa}
\usepackage{xfrac}
\usepackage[unicode,hidelinks,bookmarks=false]{hyperref}
\newcommand{\ie}{i.e.\ }
\newcommand{\cf}{cf.\ }
\newcommand{\resp}{respectively\ }
\newcommand{\Subsection}{\subsection}
\providecommand{\nobreakdash}{\nobreak\hbox{-}}
\setlength{\emergencystretch}{2em}
\multlinegap0pt
\usepackage{amssymb}
\newtheorem{lemma}{Lemma}
\title{A sharp $p$-subadditive bound for the $\ell_p$ Hausdorff distance from convex hull}
\author{Mark Meyer}
\date{Source: arXiv:2604.20387v2 (CC BY 4.0)}
\begin{document}
\maketitle

\noindent\emph{Source and licence.} A sharp $p$-subadditive bound for the $\ell_p$ Hausdorff distance from convex hull. Version arXiv:2604.20387v2 (CC BY 4.0), distributed by arXiv under the Creative Commons Attribution 4.0 International licence, http://creativecommons.org/licenses/by/4.0/, as recorded in the arXiv licence metadata for this exact version and shown on the abstract page https://arxiv.org/abs/2604.20387v2. Full licence notice: \texttt{licenses/arxiv-2604.20387-CC-BY-4.0.txt}.\par\smallskip

\noindent\textit{Editorial scope.} Counted unit: the article's lemma lemma:parallelogram\_reduction\_K\_norm (source0.tex lines 207--221). The article's complete original proof of it is delivered with it (source0.tex lines 223--261, one \texttt{proof} environment of 39 lines). No mathematics is written, repaired or re-derived: the statement and the proof are the article's own lines, in the article's own notation, and no retained line is substituted or re-flowed.  Retained uncounted as the source-local context the proof needs: lines 122--131; lines 151--157; lines 169--179.  Nothing was omitted from any retained span, so the package carries no omission record.  No retained line contains a citation, so the wrapper carries no bibliography and no bibliography entry.  No retained line contains a figure environment, an image-inclusion command or drawing code, so no image is delivered and the package declares no asset dependency.  Editorial formatting only, in addition: an AMS \texttt{amsart} preamble that restates the article's own declarations and macros used by the retained lines, namely the environments \texttt{lemma} on separate counters, together with the standard packages those lines need.  The retained environments print in this document as Lemma 1 in order of appearance, whereas the article prints the same items with the numbering of its own full document.  The complete native source of the article as distributed in the arXiv e-print for this exact version is bundled unchanged as \texttt{sources/198928/source0.tex}.
\begin{lemma}\label{lemma:opposite_hyperplane}
    Let $K\subset\mathbb{R}^n$ be a convex body that is symmetric about the origin, let $B\subset\mathbb{R}^n$ be a nonempty compact set, and let $H$ be a hyperplane that intersects $\textup{conv}(B)$. If $H_1$ and $H_2$ are the two closed half-spaces that are bounded by $H$, then the minimum
    \begin{align*}
        \inf\{\|b_1-b_2\|_K:b_1\in B\cap H_1,b_2\in B\cap H_2\}
    \end{align*}
    is attained, and every minimizing pair satisfies 
    \begin{align*}
        \|b_1-b_2\|_K\leq 2d^{(K)}(B).
    \end{align*}
\end{lemma}

\begin{lemma}\label{lemma:sum_decomposition}
    Let $L,V\subset\mathbb{R}^n$ be compact convex sets with $0\in L\cap V$ and $\textup{aff}(L)\subset \textup{aff}(V)$. Then
    \begin{align*}
        L+V=V\cup (\partial V+L),
    \end{align*}
    where $\partial V$ denotes the relative boundary of $V$.
\end{lemma}

\begin{lemma}\label{lemma:line_passing_technical_lemma}
    Let $A,B\subset\mathbb{R}^2$ be compact sets of positive dimension, and suppose that $x\in \textup{conv}(A+B)$ such that 
    \begin{align*}
        x\notin A+\textup{conv}(B),\qquad x\notin \textup{conv}(A)+B.
    \end{align*}
    Assume that $a_1,a_2$ are distinct elements of $A$, and that $x\in [a_1,a_2]+\textup{conv}(B)$. Set $u:=a_2-a_1$, and let $L:=x+\textup{span}(u)$. Then
    \begin{align*}
       C:= L\cap (\textup{conv}(B)+a_1)\neq \varnothing,\qquad C+[0,u]=L\cap (\textup{conv}(B)+[a_1,a_2]).
    \end{align*}
    Moreover, if $y\in C$, then $x\in y+[0,u]$, and $C\cap (B+\{a_1,a_2\})=\varnothing$.
\end{lemma}

\begin{lemma}\label{lemma:parallelogram_reduction_K_norm}
        Let $K\subset\mathbb{R}^2$ be a convex body that is symmetric about the origin, and let $A,B\subset\mathbb{R}^2$ be nonempty compact sets. If $x\in \textup{conv}(A+B)$ such that 
        \begin{align*}
            x\notin \textup{conv}(A)+B,\qquad x\notin A+\textup{conv}(B),
        \end{align*}
        then there exist distinct $a_1,a_2\in A$ and distinct $b_1,b_2\in B$ with
        \begin{align*}
            \|a_1-a_2\|_K\leq 2d^{(K)}(A),\qquad \|b_1-b_2\|_K\leq 2d^{(K)}(B),
        \end{align*}
        such that 
        \begin{align*}
            x\in [a_1,a_2]+[b_1,b_2].
        \end{align*}
        Moreover, the intervals $[a_1,a_2]$ and $[b_1,b_2]$ are not parallel.
\end{lemma}

\begin{proof}
    First, suppose that $\textup{dim}(A)=\textup{dim}(B)=1$. Write $x=a+b$, where $a\in \textup{conv}(A)$ and $b\in \textup{conv}(B)$. The assumption on $x$ implies that $a\notin A$ and $b\notin B$. Let $a_1,a_2$, and $b_1,b_2$ be the endpoints of the gaps of $A$ and $B$ that contain $a$ and $b$, in the sense that
    \begin{align*}
        a\in [a_1,a_2],\qquad b\in [b_1,b_2],\qquad \|a_1-a_2\|_K\leq 2d^{(K)}(A),\qquad \|b_1-b_2\|_K\leq 2d^{(K)}(B).
    \end{align*}
    We used that the midpoint of the gap $[a_1,a_2]$ has $K$-distance to $A$ exactly $\|a_2-a_1\|_K/2$. By definition, this distance cannot exceed $d^{(K)}(A)$, which implies the bound $\|a_2-a_1\|_K\leq 2d^{(K)}(A)$. Use an identical argument to get the estimate on $\|b_2-b_1\|_K$. Moving on, we have $x\in [a_1,a_2]+[b_1,b_2]$. Moreover, the intervals $[a_1,a_2]$ and $[b_1,b_2]$ are not parallel. Otherwise, interchange the roles of $A$ and $B$ if necessary, and label the endpoints so that $u:=a_2-a_1$ and $\lambda u=b_2-b_1$, with $\lambda\geq 1$. Then
    \begin{align*}
        [a_1,a_2]+[b_1,b_2]=a_1+b_1+([0,\lambda u]\cup [\lambda u,(1+\lambda)u]),
    \end{align*}
    so that either 
    \begin{align*}
        x\in a_1+[b_1,b_2],\qquad \textup{or}\qquad x\in b_2+[a_1,a_2],
    \end{align*}
    which in either case is a contradiction since $a_1+[b_1,b_2]\subset A+\textup{conv}(B)$ and $b_2+[a_1,a_2]\subset B+\textup{conv}(A)$.

    Next, suppose that at least one of the sets has dimension $2$. Interchanging and translating the sets if needed, we assume that $\textup{dim}(A)=2$ and that $0\in A\cap B$. Lemma \ref{lemma:sum_decomposition} gives
    \begin{align*}
        \textup{conv}(A)+\textup{conv}(B)=\textup{conv}(A)\cup (\partial \textup{conv}(A)+\textup{conv}(B)).
    \end{align*}
    Using the above representation and the assumption on $x$, write $x=a+b$, where 
    \begin{align*}
        a\in \partial\textup{conv}(A)\backslash A,\qquad b\in \textup{conv}(B).
    \end{align*}
    If $H$ is a supporting line of $\textup{conv}(A)$ that passes through $a$, then $a\in \textup{conv}(A\cap H)$. Since $A\cap H$ is a compact subset of a line and $a\notin A$, there exist distinct $\tilde{a}_1,\tilde{a}_2\in A\cap H$ with $a\in [\tilde{a}_1,\tilde{a}_2]$, and so $x\in [\tilde{a}_1,\tilde{a}_2]+\textup{conv}(B)$.

    Let $L$ be the line that passes through $x$ and is parallel to the interval $[\tilde{a}_1,\tilde{a}_2]$. By Lemmas \ref{lemma:opposite_hyperplane} and \ref{lemma:line_passing_technical_lemma}, the set $C:=L\cap (\textup{conv}(B)+\tilde{a}_1)$ is nonempty, and there exist $b_1,b_2\in B$ such that $b_1+\tilde{a}_1$ and $b_2+\tilde{a}_1$ lie on opposite sides of $L$ and satisfy
    \begin{align*}
        \|b_1-b_2\|_K\leq 2d^{(K)}(B).
    \end{align*}
    Now, the interval $\tilde{a}_1+[b_1,b_2]$ intersects $L$ at a point $y$. Then $y\in C$, and by Lemma \ref{lemma:line_passing_technical_lemma}, we have $x\in y+[0,\tilde{a}_2-\tilde{a}_1]$, and so
    \begin{align}\label{eq:halfway_done}
        x\in [\tilde{a}_1,\tilde{a}_2]+[b_1,b_2].
    \end{align}
    Repeat the same argument with the roles of $A$ and $B$ reversed, starting with \eqref{eq:halfway_done} to find $a_1,a_2\in A$ such that
    \begin{align*}
        \|a_1-a_2\|_K\leq 2d^{(K)}(A),\qquad x\in [a_1,a_2]+[b_1,b_2].
    \end{align*}
    We use the same argument from the beginning of this proof to show that the intervals must be nonparallel.
\end{proof}

\end{document}
